\scanpage{226}
\mspage{172}
% archive sheet number: 225

% MS lines 1–5, reference and notation
Pour la notion de \deleted{\unsure{f}} morphismes $W$-propre, $W$-lisse, $W$-\unsure{parfaits|parfait|parfaites} dans $\mathrm{Cat}$, je renvoie à $\overline{\underline{\mathrm{XII}}}$, partie $\overline{\underline{\mathrm{IV}}}$ (p. 22-52), et les notations $\underline{\mathrm{Liss}}_W S$, $\underline{\mathrm{Prop}}_W S$, $\underline{\mathrm{Parf}}_W S$ \unsure{etc.|e} (sous-catégories de $\mathrm{Cat}_{/S}$), et $\mathrm{Parf}_W\subset\mathrm{Fl}(\mathrm{Cat})$. Dans l'énoncé de $W(4\,\mathrm b')$, on a posé

% MS lines 6–8, displayed definition
\[
\begin{aligned}
\underline{\mathrm{Parf}}_{\unsure{\omega|\infty}}S
 &=\underline{\mathrm{Parf}}_{W_\omega}S,\quad \text{où }W_\omega\subset\mathrm{Fl}(\mathrm{Cat})\text{ est la}\\
 &\hspace{4em}\text{plus petite partie de }\mathrm{Fl}(\mathrm{Cat})\\
 &\hspace{4em}\text{satisfaisant }W(1,2\,\mathrm{bis},3).
\end{aligned}
\]

% MS line 9
Je note de même

% MS lines 10–11, displayed definition
\[
\begin{aligned}
W'_\omega\subset\mathrm{Fl}(\mathrm{Cat}),\quad &\text{plus petite partie de }\mathrm{Fl}(\mathrm{Cat})\\
 &\text{satisfaisant }W(1,2,3,4\,\mathrm b).
\end{aligned}
\]

% MS lines 12–15, inclusion and replacement
L'hypothèse $W\supset W_\infty$ que nous avions fini par imposer sur $W$, et qui \deleted{\unsure{paraît}} \inserted{est} artificielle \unsure{au} \deleted{\unsure{vu}} \inserted{cours} \unsure{de la réflexion|des réflexions} de $\overline{\underline{\mathrm{XII}}}$, a avantage à être remplacée par

% MS line 16, inclusion
\[
 W\supset W'_\omega\quad \text{(ou par }W\supset\overline{W'_\omega},\ \text{où }\overline{W'_\omega}\text{ est le saturé fort de }W'_\omega\text{).}
\]

% MS left margin of middle paragraph, upward writing
% slanted
\margin{$\overline{W'_\omega}$ est le plus petit \unsure{sous-ensemble} \uline{fortement} \uline{saturé} de $\mathrm{Fl}(\mathrm{Cat})$ \unsure{satisfaisant} $W(4\,\mathrm b)$, et \deleted{\unsure{dans}} il \unsure{satisfait} à la condition $W'(\unsure{4.|4})$, variante plus forte \unsure{des deux axiomes|des axiomes} de \unsure{la forme} (4).}

% MS lines 17–21, optional condition
Mais, nous ne l'imposons pas$^*$ comme condition liminaire, au même titre que $W(1,2,3)$, pas plus que $W(4\,\mathrm b)\unsure{^*|^{\prime}}$. Le rôle que j'ai fait jouer à $W_\infty$ dans les semaines passées, semble revenir plus naturellement à $W'_\omega$ ou $\overline{W'_\omega}$.

% MS lines 22–last, Der 3
Il en est ainsi notamment \deleted{\unsure{aussi}} \inserted{aussi,} il me semble, dans l'énoncé des conditions du type Der 3 sur un dérivateur ($\overline{\mathrm I}$, p. 111-112), qu'il faudrait revoir dans cette lumière. Mais tout dépend si \inserted{(comme j'en ai l'impression)} on peut faire vraiment \deleted{\unsure{une}}\inserted{de} l'algèbre homotopique substantielle, avec notamment les suites exactes longues de fibration et de suspension, pour $\mathrm{HOT}_W$, avec des $W$ qui sont \uline{strictement} contenus dans $W_\infty$. Avant de m'en être assuré, il est trop tôt encore pour me prononcer sur la ``bonne'' formulation de Der 3,
